\documentclass[11pt,letterpaper]{article}
\usepackage[utf8]{inputenc}
\usepackage[spanish]{babel}
\usepackage{amsmath,amssymb,amsfonts}
\usepackage{geometry}
\geometry{top=2.5cm, bottom=2.5cm, left=3cm, right=3cm}
\usepackage{hyperref}
\usepackage{booktabs}
\usepackage{color}

\definecolor{darkblue}{rgb}{0.0, 0.0, 0.55}

\hypersetup{
    colorlinks=true,
    linkcolor=darkblue,
    filecolor=magenta,      
    urlcolor=darkblue,
    citecolor=darkblue,
}

\title{\textbf{Compendio Avanzado sobre Teoría de Grafos y Algoritmos de Enrutamiento: Desde los Fundamentos Clásicos hasta el Estado del Arte Topológico}}
\author{\textbf{Reporte de Investigación Científica y Tecnológica}}
\date{\today}

\begin{document}

\maketitle

\begin{abstract}
Este reporte presenta una revisión exhaustiva y rigurosa de la teoría de grafos y los algoritmos de enrutamiento optimizado. El análisis abarca desde las definiciones topológicas elementales y las representaciones matriciales computacionales hasta las metodologías contemporáneas de aprendizaje profundo en grafos (Graph Neural Networks). Se exploran en profundidad los algoritmos clásicos de caminos mínimos, flujos en redes, análisis de centralidad y optimización combinatoria, proporcionando formulaciones matemáticas explícitas, análisis de complejidad algorítmica y comparaciones computacionales. El objetivo fundamental es consolidar un marco de referencia analítico para el diseño, simulación y optimización de redes complejas e infraestructuras industriales.
\end{abstract}

\newpage
\tableofcontents
\newpage

\section{Introducción y Fundamentos Topológicos}
La teoría de grafos constituye el pilar formal para el modelado, análisis y optimización de sistemas interconectados. Estructuras como redes de distribución eléctrica, sistemas de automatización industrial, redes logísticas y sistemas de transporte pueden ser completamente caracterizadas mediante estas herramientas matemáticas discretas.

\subsection{Definiciones Formales}
Formalmente, un grafo se define como un par ordenado $G = (V, E)$, donde:
\begin{itemize}
    \item $V$ es un conjunto no vacío de \textbf{vértices} o \textbf{nodos}, que representan las entidades del sistema (e.g., subestaciones, motores, servidores, nodos de control).
    \item $E$ es un conjunto de \textbf{aristas}, \textbf{arcos} o \textbf{enlaces}, que representan las relaciones o conexiones físicas/lógicas entre dichas entidades, de tal manera que $E \subseteq \{\{u, v\} \mid u, v \in V\}$.
\end{itemize}

\subsection{Clasificación de Estructuras}
Dependiendo de la naturaleza de las interconexiones, los grafos se clasifican estructuralmente en:
\begin{enumerate}
    \item \textbf{Grafos No Dirigidos:} Las aristas no poseen una orientación intrínseca; la relación entre $u$ y $v$ es simétrica ($(u,v) \equiv (v,u)$).
    \item \textbf{Grafos Dirigidos (Digrafos):} Las aristas poseen una dirección definida, representadas como pares ordenados $(u, v)$ donde el flujo u operación se origina en $u$ (nodo cola) y termina en $v$ (nodo cabeza).
    \item \textbf{Grafos Ponderados:} Cada arista $e \in E$ tiene asignada una función de costo o peso $w(e) \in \mathbb{R}$. En entornos de ingeniería, este peso representa magnitudes físicas o financieras medibles, tales como distancia en metros, caída de tensión, resistencia óhmica, latencia de red o costos de construcción.
\end{enumerate}

\subsection{Representaciones Computacionales}
La eficiencia de los algoritmos de enrutamiento depende críticamente de la estructura de datos subyacente utilizada para almacenar el grafo en memoria:

\subsubsection{Matriz de Adyacencia}
Para un grafo con $|V| = n$, la matriz de adyacencia $A$ es una matriz de dimensiones $n \times n$ donde el elemento $A_{ij}$ se define como:
\begin{equation}
A_{ij} = \begin{cases} 
w(v_i, v_j) & \text{si } (v_i, v_j) \in E \\
0 \text{ o } \infty & \text{si } (v_i, v_j) \notin E 
\end{cases}
\end{equation}
Esta representación requiere un espacio en memoria de $\mathcal{O}(V^2)$, lo que la hace altamente ineficiente para \textbf{grafos dispersos} ($|E| \ll |V|^2$). Sin embargo, permite consultas de existencia de aristas en tiempo constante $\mathcal{O}(1)$.

\subsubsection{Lista de Adyacencia}
Asocia a cada vértice $v \in V$ una lista de sus vecinos directos. Su consumo espacial es de $\mathcal{O}(V + E)$, siendo la opción óptima para la gran mayoría de los algoritmos de exploración topológica.

\subsubsection{Matrices Dispersas (Sparse Matrices)}
En la computación científica de alto rendimiento, los grafos masivos se representan mediante esquemas de matrices dispersas como Compressed Sparse Row (CSR) o Compressed Sparse Column (CSC), optimizando el uso de memoria caché y permitiendo operaciones de álgebra lineal distribuida.

\section{Algoritmos Clásicos de Caminos Mínimos y Exploración}
La determinación de la trayectoria óptima entre nodos es el núcleo operativo del enrutamiento. 

\subsection{Algoritmo de Dijkstra}
Diseñado por Edsger Dijkstra, es el algoritmo de referencia para encontrar los caminos mínimos desde un único nodo origen hacia todos los demás vértices en un grafo ponderado con pesos no negativos ($w(e) \ge 0$). La condición de relajación para una arista $(u, v)$ con peso $w(u,v)$ se expresa como:
\begin{equation}
\text{si } d(v) > d(u) + w(u, v) \implies d(v) = d(u) + w(u, v)
\end{equation}
Utilizando una cola de prioridad, la complejidad temporal del algoritmo se optimiza a $\mathcal{O}((V + E) \log V)$.

\subsection{Algoritmo A* (A-Star)}
Introducido por Hart, Nilsson y Raphael, el algoritmo A* introduce una función de evaluación $f(n) = g(n) + h(n)$ para determinar el orden de exploración, donde $h(n)$ es una heurística que estima el costo restante más bajo desde el nodo $n$ hasta el destino final.

\subsection{Algoritmo de Bellman-Ford}
A diferencia de Dijkstra, el algoritmo de Bellman-Ford puede computar caminos mínimos en grafos que contienen aristas con pesos negativos. Su funcionamiento se basa en una aproximación por programación dinámica. Es vital para \textbf{detectar ciclos de costo negativo}. Su complejidad temporal es de $\mathcal{O}(V \cdot E)$.

\section{Árboles de Expansión Mínima y Flujo en Redes}

\subsection{Árbol de Expansión Mínima (MST)}
Esquemas centrales para el diseño eficiente de infraestructuras físicas minimizando el uso de materiales y conexiones.
\begin{itemize}
    \item \textbf{Algoritmo de Prim:} Complejidad $\mathcal{O}(E + V \log V)$. Basado en la expansión desde nodos adyacentes.
    \item \textbf{Algoritmo de Kruskal:} Complejidad $\mathcal{O}(E \log E)$. Enfocado en el ordenamiento de arcos de menor costo.
\end{itemize}

\subsection{Algoritmos de Flujo Máximo}
El \textbf{Teorema del Flujo Máximo - Corte Mínimo} rige la determinación de capacidades límite operativas en sistemas. Algoritmos destacados incluyen \textbf{Ford-Fulkerson} y \textbf{Edmonds-Karp} para el modelado de caudales y dimensionamiento de enlaces de red.

\section{Análisis de Centralidad y Optimización Combinatoria}
\subsection{Métricas de Centralidad}
\begin{itemize}
    \item \textbf{Centralidad de Intermediación (Betweenness):} Cuantifica cuántas veces un componente es un eslabón crítico para las rutas más cortas de la red.
    \item \textbf{PageRank:} Evalúa la dominancia estructural a partir de matrices probabilísticas de transición.
\end{itemize}

\section{El Estado del Arte: Aprendizaje Profundo en Grafos (GNN)}
El cruce entre IA topológica y matrices estocásticas se define mediante redes neuronales de paso de mensajes (Message Passing).
\begin{equation}
H^{(l+1)} = \sigma \left( \tilde{D}^{-\frac{1}{2}} \tilde{A} \tilde{D}^{-\frac{1}{2}} H^{(l)} W^{(l)} \right)
\end{equation}
Las Redes Convolucionales de Grafos (GCN) representan el estándar moderno para mantenimiento predictivo y modelado hiper-dimensional de topologías técnicas complejas.

\begin{thebibliography}{99}

\bibitem{bondy2008}
J.A. Bondy y U.S.R. Murty.
\emph{Graph Theory}.
Graduate Texts in Mathematics, Springer, 2008.

\bibitem{cormen2022}
T.H. Cormen, C.E. Leiserson, R.L. Rivest y C. Stein.
\emph{Introduction to Algorithms}.
MIT Press, 4ta Edición, 2022.

\bibitem{kipf2017}
T.N. Kipf y M. Welling.
Semi-supervised classification with graph convolutional networks.
\emph{ICLR}, 2017.

\bibitem{ford2010}
L.R. Ford Jr. y D.R. Fulkerson.
\emph{Flows in Networks}.
Princeton University Press, 2010.

\end{thebibliography}

\end{document}
